\textbf{Two evaluation defects and their correction.} We report both because
each is easy to reproduce inadvertently and neither is visible in the resulting
scores.

The first concerns failed judge calls: when a call errored, the pipeline
persisted a score record containing five zeros and continued, and because
subsequent runs skip cases for which a score record already exists, those
failures became permanent. In our initial pass this silently invalidated 231 of
the test-split scores, 211 of which belonged to a single model, so that the
model's reported row was computed from 309 of 550 cases while the accompanying
table indicated 520. The judge now declines to persist a failed score, leaving
the case pending for a subsequent run.

The second concerns how reasoning traces are delivered. Our post-trained model
was served through vLLM without a reasoning parser, and because Qwen chat
templates pre-fill \texttt{<think>} in the generation prompt, the model emits its
reasoning, a closing \texttt{</think>}, and then its answer within a single
field, whereas every hosted API returns reasoning in a separate field. Our model
was consequently graded on approximately 4,000 characters of unfiltered reasoning
prepended to its answer while all other systems were graded on the answer alone;
correcting this raised its test-split score from 3.81 to 4.17 and moved it from
third position to second. The general point is that comparisons between locally
served reasoning models and hosted APIs require verifying the contents of the
string being graded.

All results reported in this paper follow both corrections: every score in the study is
produced by GPT-5.4, and no score record contains a judge failure.
